% Section 5 — Results on synthetic microservices telemetry (OB)
% Page budget: ~3 pages in acmart sigconf two-column (~2500-2800 words).
% Source-of-truth (every number): the per-section SUMMARY.md files under
%   results/ob/{3.5..4.14}/SUMMARY.md and
%   DOCS/docs8/PAPER-FINDINGS.md.
%
% Editorial discipline:
%   - Every measured number cites a source-of-truth file in a footnote
%     or in-line. No fabrication; if a value cannot be located in the
%     source, it does not appear.
%   - Cross-dataset numbers (OTel/WoL) live in §6. §5 is OB-deep with
%     brief cross-confirmations where they strengthen the within-dataset
%     story.

\section{Results}
\label{sec:results}

All numbers below are computed on the held-out test partition of \dsob{} (\S\ref{sec:evaluation:datasets:ob}) unless noted. Headline metrics carry $95\%$ percentile confidence intervals from $1{,}000$-resample paired bootstrap at seed $42$; pairwise deltas additionally carry \emph{paired-delta} confidence intervals as described in \S\ref{sec:evaluation:stats}.

\subsection{Headline performance}
\label{sec:results:headline}

Table~\ref{tab:headline-ob} reports the agent's headline metrics on the $1{,}008$-window \dsob{} test split (of which $331$ are evaluable for retrieval --- they carry at least one time-visible gold ticket).

\begin{table}[ht]
  \centering
  \caption{Headline metrics on \dsob{}. $95\%$ percentile CIs from $1{,}000$-resample bootstrap at seed $42$.}
  \label{tab:headline-ob}
  \small
  \begin{tabular}{lrr}
    \toprule
    Metric & Point & $95\%$ CI \\
    \midrule
    Hit@1                       & $0.6888$ & $[0.6385, 0.7405]$ \\
    Hit@5                       & $0.7583$ & $[0.7085, 0.8029]$ \\
    Hit@10                      & $0.7583$ & $[0.7085, 0.8029]$ \\
    MRR                         & $0.7138$ & $[0.6637, 0.7604]$ \\
    Triage accuracy             & $0.8343$ & $[0.8095, 0.8591]$ \\
    Novel recall / precision    & $0.5361 / 0.1128$ & --- \\
    Pages-per-incident          & $1.000$  & --- \\
    Suppressions fired          & $20$     & --- \\
    Distinct plan identifiers   & $4 / 6$  & --- \\
    \bottomrule
  \end{tabular}
\end{table}

Three numbers in this table are load-bearing for the rest of the paper. First, Hit@5 at $0.7583$ is the retrieval ceiling we will hold constant while we vary cost in \S\ref{sec:results:cost} and \S\ref{sec:results:pareto}: every optimisation the agent permits has to preserve this value or it will not appear in the production recommendation. Second, triage accuracy at $0.8343$ is the metric we will see degrade significantly when we remove the numeric-feature classifier in \S\ref{sec:results:ablation}. Third, pages-per-incident at $1.000$ across all twenty test-set suppression events is the operational claim we will replicate on the other two datasets in \S\ref{sec:cross-app}. Novelty precision is low ($0.1128$); this is a known blind spot whose mechanism we localise in \S\ref{sec:results:failures}.

\subsection{Adaptive selection cuts cost without quality loss}
\label{sec:results:cost}

The headline cost claim is computed by comparing the agent's actual cost on the \dsob{} test split against an always-on cascade counterfactual in which every diagnostic primitive runs on every window (see \S\ref{sec:evaluation:baselines}). Costs are sourced from a per-skill table whose assumptions are documented in Appendix~A; the language-model verifier dominates the dollar axis at $\$0.0005$ per call and $1{,}969$ tokens.

Table~\ref{tab:cost-vs-cascade} reports cost in two configurations. The \emph{verifier-off} configuration is the cheap-smoke default (the verifier is not registered, mostly because the smoke does not require the LM Studio service to be reachable). The \emph{verifier-on} configuration is the production-realistic one in which the always-on baseline does include the language-model verifier; this is the configuration the paper recommends as the headline cost result, because the verifier is the configuration that meaningfully matters in deployment.

\begin{table}[ht]
  \centering
  \caption{Cost on \dsob{}: agent versus the always-on cascade counterfactual on the same $1{,}008$ windows.}
  \label{tab:cost-vs-cascade}
  \small
  \setlength{\tabcolsep}{4pt}
  \begin{tabular}{lrrr}
    \toprule
    Configuration & Resource & Agent & Cascade \\
    \midrule
    Verifier-off  & Wall time (s)       & $61.2$    & $101.9$   \\
                  & LLM tokens          & $71{,}600$ & $121{,}600$ \\
                  & USD-equivalent      & $\$0.018$  & $\$0.030$ \\
    \cmidrule(lr){1-4}
    \textbf{Verifier-on} & Wall time (s) & $1{,}524$ & $4{,}783$ \\
                  & LLM tokens          & $258{,}655$ & $720{,}176$ \\
                  & USD-equivalent      & $\$0.065$  & $\$0.182$  \\
    \midrule
    \multicolumn{4}{l}{\textbf{Savings (verifier-on, production-realistic): $68.1\%$ wall, $64.1\%$ tokens, $64.1\%$ USD.}} \\
    \bottomrule
  \end{tabular}
\end{table}

The savings split across two mechanisms. The first, capability-drop, removes a primitive from the plan when its required evidence is not in the bundle: on \dsob{}, the structural and knowledge-graph retrievers are dropped because the dataset does not surface a knowledge-graph capability flag (which is consistent with the agent's design \S\ref{sec:approach:gating} and disclosed honestly in the per-skill breakdown). The second mechanism, gate-skip, drops a primitive at runtime when the cheap path is sufficiently confident. The gate-skip rates per primitive are reported in Table~\ref{tab:per-skill-savings}.

\begin{table}[ht]
  \centering
  \caption{Per-primitive gate-save rates on \dsob{} (verifier-on). A high save rate means the cheap path was confident on most active-fault windows for that primitive.}
  \label{tab:per-skill-savings}
  \small
  \begin{tabular}{lrrr}
    \toprule
    Primitive & Invoked & Gated out & Save rate \\
    \midrule
    Hybrid retriever (rule)   & $95$  & $209$ & $68.8\%$ \\
    Language-model verifier   & $95$  & $209$ & $68.8\%$ \\
    Reformulator              & $179$ & $125$ & $41.1\%$ \\
    Each of the 4 ReAct tools & $179$ & $125$ & $41.1\%$ \\
    Reranker                  & $179$ & $125$ & $41.1\%$ \\
    Retrieval composer        & $678$ & $31$  & $4.4\%$  \\
    Dense retriever, triage classifier, composers & $1008$ & $0$ & $0.0\%$ \\
    \bottomrule
  \end{tabular}
\end{table}

The structural pattern is the one the agent's design predicts: the most expensive primitives have the highest save rates. The language-model verifier and the hybrid retriever, both gated by the same cheap-path-confident condition, skip on two thirds of active-fault windows. The four evidence-gathering tools and the reformulator share a lower-consensus gate and skip on roughly forty percent. The cheap-path primitives (the dense retriever, the numeric classifier, the composers) run on every window because they are the cheap path. Hit@K is unchanged across this entire skip pattern: the gated primitives are not contributing to the rank-one ranking on the windows where they fire, which we confirm directly in \S\ref{sec:results:ablation}.

\subsection{The cost-quality Pareto curve}
\label{sec:results:pareto}

A single threshold parameter controls how confident the cheap path must be before the gate closes. Sweeping this threshold over $\{0.50, 0.60, 0.70, 0.80, 0.90, 0.95\}$ traces a Pareto curve over (Hit@$K$, cost). Table~\ref{tab:pareto-ob} reports the curve under the verifier-on configuration.

\begin{table}[ht]
  \centering
  \caption{Pareto sweep over the cheap-path confidence threshold on \dsob{} (verifier-on). Hit@K is invariant; only cost varies.}
  \label{tab:pareto-ob}
  \small
  \begin{tabular}{lrrrr}
    \toprule
    Threshold & Hit@1 & Hit@5 & USD & Wall saved \\
    \midrule
    $0.50$              & $0.6888$ & $0.7583$ & $\$0.0589$ & $72.3\%$ \\
    $0.60$              & $0.6888$ & $0.7583$ & $\$0.0599$ & $71.7\%$ \\
    $0.70$              & $0.6888$ & $0.7583$ & $\$0.0614$ & $70.7\%$ \\
    $0.80$              & $0.6888$ & $0.7583$ & $\$0.0624$ & $70.1\%$ \\
    $0.90$ (default)    & $0.6888$ & $0.7583$ & $\$0.0654$ & $68.1\%$ \\
    $0.95$              & $0.6888$ & $0.7583$ & $\$0.0679$ & $66.5\%$ \\
    \bottomrule
  \end{tabular}
\end{table}

The Pareto front is \emph{flat} in Hit@K across the entire operating range. Every retrieval metric is identical to four decimal places for every threshold; only USD cost varies. The agent's design predicts this: the cascade-derived primitives that get gated as the threshold drops do not change the rank-one ranking on \dsob{} (see \S\ref{sec:results:ablation}), so reducing the threshold makes the gate close more often without sacrificing retrieval quality. The production recommendation is therefore unambiguous: drop the threshold from the documented default of $0.90$ to $0.50$ and recover an additional four percentage points of cost savings at zero retrieval cost. We discuss the cross-dataset generalisability of this recommendation in \S\ref{sec:cross-app}.

\subsection{What carries the lift: skill and tool ablation}
\label{sec:results:ablation}

The cost-savings result raises an obvious question: \emph{which} primitives are load-bearing and which are decoration? We answer it with three ablation grids: one over individual cascade-derived primitives, one over the four-tool on-demand evidence-gathering subset, and one over capability flags.

\subsubsection{Skill ablation}

The skill-ablation grid disables one primitive (or one named group) at a time and re-runs the agent on the same $1{,}008$ test windows. Of ten ablation cells tested, only two produce a non-zero paired-delta against the baseline. Table~\ref{tab:skill-paired} reports the paired-delta confidence intervals for the four cells that move at least one metric.

\begin{table}[ht]
  \centering
  \caption{Skill-ablation paired-delta CIs on \dsob{} (1{,}000 paired resamples). Cells not shown produce exactly $0$ per-window delta on every metric; their CI collapses to $[0, 0]$ with $p=1.000$.}
  \label{tab:skill-paired}
  \small
  \begin{tabular}{llrrr}
    \toprule
    Cell & Metric & $\Delta$ & $95\%$ CI & $p$ \\
    \midrule
    \texttt{no\_react}           & Hit@1     & $-0.0121$ & $[-0.0302, +0.0060]$ & $0.262$ \\
    \texttt{no\_react}           & MRR       & $-0.0078$ & $[-0.0199, +0.0033]$ & $0.176$ \\
    \texttt{no\_triage\_numeric} & Triage    & $-0.117$  & $[-0.139, -0.097]$   & $<10^{-4}$ \\
    \texttt{dense\_only}         & Triage    & $-0.117$  & $[-0.139, -0.097]$   & $<10^{-4}$ \\
    \bottomrule
  \end{tabular}
\end{table}

Two findings drop out. First, the numeric-feature classifier is statistically essential for triage accuracy: removing it (either directly or via the \texttt{dense\_only} reduction) drops triage accuracy by $11.7$ absolute points with $p < 10^{-4}$. The effect is large, sharp, and one-sided. Second, the on-demand evidence-gathering loop's positive contribution is \emph{not} statistically significant: the point estimate of $-0.0121$ Hit@1 for the \texttt{no\_react} cell straddles zero in its $95\%$ paired interval, with $p = 0.262$. We cannot reject the null hypothesis that the loop has no effect on Hit@1.

Every other skill ablation produces exactly zero per-window delta across all metrics. Removing the structural retriever, the knowledge-graph retriever, the log-sequence retriever, the hybrid-fusion variants, the reformulator, the language-model verifier, or the novelty composer leaves every test window's retrieval rank unchanged. The cascade's retrieval ranking on \dsob{} is anchored by the dense retriever, and the other primitives serve as backup signals on the few windows where the dense retriever is genuinely uncertain.

\subsubsection{Tool ablation}

The on-demand evidence-gathering loop comprises four tools (Kubernetes events, extended trace, Prometheus snapshot, peer-incident retrieval) plus a reranker. The four tools yield $2^{4} = 16$ activation subsets. Table~\ref{tab:tool-paired} reports the paired-delta CIs against the empty (no-tools) baseline for the most informative cells; the full table appears in the supplementary material.

\begin{table}[ht]
  \centering
  \caption{Tool-subset paired-delta CIs on \dsob{} (against the empty no-tools baseline). The single statistically-significant finding is a \emph{negative} one: the extended-trace tool harms Hit@1.}
  \label{tab:tool-paired}
  \small
  \begin{tabular}{lrrr}
    \toprule
    Subset & $\Delta$ Hit@1 & $95\%$ CI & $p$ \\
    \midrule
    \textbf{Extended-trace alone}      & $\mathbf{-0.0181}$ & $\mathbf{[-0.0332, -0.0060]}$ & $\mathbf{0.002}$ \\
    Extended-trace $+$ metrics         & $-0.0181$ & $[-0.0332, -0.0060]$ & $0.002$ \\
    Events $+$ extended-trace          & $-0.0121$ & $[-0.0272,  0.0000]$ & $0.104$ \\
    Peers alone (best by point estimate) & $+0.0151$ & $[-0.0060, +0.0363]$ & $0.180$ \\
    All four tools (full ReAct config) & $+0.0121$ & $[-0.0060, +0.0302]$ & $0.262$ \\
    Events alone                       & $+0.0030$ & $[-0.0091, +0.0181]$ & $0.832$ \\
    Metrics alone                      & $\;0$     & $[0, 0]$             & $1.000$ \\
    \bottomrule
  \end{tabular}
\end{table}

The single statistically-significant ReAct-related finding on \dsob{} is a negative one: invoking the extended-trace tool in isolation (or with the metrics tool) costs $1.8$ absolute Hit@1 percentage points at $p = 0.002$, with a paired CI strictly below zero. The mechanism is intelligible: the trace tool surfaces a list of services seen on the call path during the window, and on \dsob{}'s call topology those service names match memory-side tickets across multiple unrelated scenario families, polluting the reranker's token bag with cross-family noise. The peer-incident tool, by contrast, surfaces same-family memory-side prose whose tokens are family-specific, and its point estimate is positive; but the positive estimate ($+0.0151$ Hit@1) does not separate from noise ($p = 0.180$), and we report it as suggestive rather than conclusive.

The honest paper claim that emerges from Tables~\ref{tab:skill-paired} and~\ref{tab:tool-paired} is that the on-demand evidence-gathering loop's framework is fully implemented, fully instrumented, fully replayable, and---on this dataset, at this sample size---does not deliver a statistically distinguishable Hit@1 lift. One of its tools is significantly harmful and should be disabled in production. Per-family tool selection (\S\ref{sec:discussion}) is the natural follow-up that this finding motivates.

\subsection{The budget-bounded retrieval curve}
\label{sec:results:budget}

A complementary view of the four-tool subset is the budget-bounded retrieval curve, in which we cap the number of tools that may be invoked per window at $N \in \{0, 1, 2, 3, 4\}$ and let the controller's branch list pick which $N$ to call. Table~\ref{tab:budget-ob} reports the curve.

\begin{table}[ht]
  \centering
  \caption{Budget-bounded Hit@K on \dsob{}. The curve is non-monotone in the number of tools called.}
  \label{tab:budget-ob}
  \small
  \begin{tabular}{rrrrr}
    \toprule
    $N$ & Hit@1 & Hit@5 & MRR & Triage acc \\
    \midrule
    $0$ & $0.6767$ & $0.7583$ & $0.7060$ & $0.8413$ \\
    $1$ & $0.6798$ & $0.7583$ & $0.7080$ & $0.8413$ \\
    $2$ & $\mathbf{0.6647}$ & $0.7583$ & $0.6970$ & $0.8413$ \\
    $3$ & $\mathbf{0.6647}$ & $0.7583$ & $0.6970$ & $0.8413$ \\
    $4$ & $0.6888$ & $0.7583$ & $0.7138$ & $0.8413$ \\
    \bottomrule
  \end{tabular}
\end{table}

The budget curve is the strongest single illustration of the section's headline: the original ``more tools is better'' hypothesis is rejected. Adding the second and third tools on top of the first costs $1.2$ Hit@1 points relative to no tools at all; only the fourth tool (peer-incident retrieval) recovers and overshoots. The mechanism is the same as Table~\ref{tab:tool-paired}: tool two is the extended-trace tool, which surfaces cross-family service-name tokens; tool three is the Prometheus snapshot, whose synthesised symptom tokens reinforce the noise; only tool four's same-family peer prose overpowers the noise. We discuss the implications for per-family tool selection in \S\ref{sec:discussion}.

\subsection{Capability-mask robustness}
\label{sec:results:capmask}

The capability-mask sweep tests the agent's graceful degradation: we strip named capability flags from the inspector's output before the controller emits a plan, then re-run the agent under the counterfactual condition. Of the eight masks tested, two move metrics. Removing the numeric-feature capability collapses triage accuracy by $11.7$ absolute points with $p < 10^{-4}$, replicating the skill-ablation finding through a different mechanism (the masked capability auto-drops the numeric classifier rather than removing it from the registry directly). Removing the trace-summary capability \emph{improves} Hit@1 by $0.30$ points (CI $[0.00, +0.91]$, $p = 0.65$, not significant)---the same negative-trace finding from a third angle. No other mask produces non-zero per-window deltas. The agent runs end-to-end on every mask; no crashes, no NaN outputs, and the controller's plan-identifier count stays at four across the entire grid.

The triple-validation pattern (skill ablation, tool ablation, capability mask) on the numeric-feature finding is what gives us high confidence in the load-bearing claim: triage accuracy is carried by the numeric-feature classifier with $p < 10^{-4}$ from three independent mechanisms, and no other component reaches significance.

\subsection{Against the naive baseline}
\label{sec:results:bm25}

A reviewer's first question for any retrieval paper is whether the system beats a naive keyword baseline. Table~\ref{tab:bm25} reports the BM25-only retrieval against the agent on the agent-evaluable subset of each test split. The BM25 entries are computed on the identical window set as the agent's so the comparison is apples-to-apples.

\begin{table}[ht]
  \centering
  \caption{Naive-baseline comparison: BM25-only retrieval versus the agent, on the agent-evaluable subset of each test split.}
  \label{tab:bm25}
  \small
  \begin{tabular}{lrrrrr}
    \toprule
    Dataset & $n_{\text{eval}}$ & BM25 Hit@1 & BM25 Hit@5 & Agent Hit@5 & Lift \\
    \midrule
    \dsob{}    & $331$ & $0.0514$ & $0.0876$ & $0.7583$ & $\mathbf{8.66\times}$ \\
    \dsotel{}  & $119$ & $0.3866$ & $0.5630$ & $0.7563$ & $1.34\times$ \\
    \dswol{}   & $55$  & $0.0000$ & $0.6000$ & $0.8364$ & $1.39\times$ \\
    \bottomrule
  \end{tabular}
\end{table}

On the synthetic dataset BM25 collapses to Hit@5 below $9\%$ because the humanised Jira tickets use vocabulary that does not match the telemetry-side evidence text; the agent's $0.7583$ is $8.66\times$ better. On the real-Apache-Jira dataset BM25 retrieves the gold ticket in the top-five $60\%$ of the time but \emph{never} as the top-one ($\text{Hit@1} = 0$), because keyword overlap alone is insufficient to disambiguate among same-project tickets; the agent's $0.7818$ Hit@1 makes the difference between a system that always asks the engineer to pick and one that puts the right answer first $78\%$ of the time. The cross-dataset replication is the subject of \S\ref{sec:cross-app}.

\subsection{Operational metrics and failure modes}
\label{sec:results:failures}

We close the OB results with the operational metrics and the failure-mode catalog.

\paragraphhead{Plan diversity (RQ-A1).} The controller emits four distinct plan identifiers across the $1{,}008$ \dsob{} test windows. The six branches defined in \S\ref{sec:approach:plan} are all implemented; the two that do not fire are the state-suppression branch (which requires a consecutive ticket-worthy sequence on the same scenario family without recovery, a pattern the locked test split does not surface frequently) and the default branch (which catches unrecognised window types, of which the split has none). The plan-identifier count is consistent with a controller that branches on real input rather than emitting a single fixed pipeline.

\paragraphhead{Page suppression (RQ-A4).} The cross-window state layer fires twenty suppression events on the test split, downgrading repeated pages to attached-to-existing-incident status. The resulting pages-per-incident ratio is $1.000$ ($358$ pages over $358$ incidents). The metric is universal across the three datasets, as \S\ref{sec:cross-app} confirms.

\paragraphhead{Tool-use failure modes (RQ-D6).} Of $716$ tool invocations on the active-fault gated subset, $670$ ($93.6\%$) return a usable signal, $46$ ($6.4\%$) return an empty result, and the remaining three failure modes (hallucinated tool name, looping repeated call, budget exhausted) each fire on zero windows. The empty-result cases all concentrate on the peer-incident tool applied to the loadgenerator scenario family, where the post-own-incident-exclusion peer set is empty by construction. The three never-firing failure modes are defences for the v2 language-model-emitted version of the loop (\S\ref{sec:approach:tooluse}); v1's rule-based controller cannot reach them by construction.

\paragraphhead{Failure-category distribution (RQ-D5).} Categorising every test-set decision into one of eight primary failure or success buckets, the distribution is dominated by two cells: $34.6\%$ of windows are correctly dismissed as noise, and $32.5\%$ are flagged as novel when in fact a memory-side gold exists (the false-novel category). The false-novel rate is the agent's largest blind spot. We localise its mechanism in \S\ref{sec:discussion}: the false-novel signal is driven by the language-model verifier's overly-eager novelty flag, which is structurally absent on datasets where the verifier is disabled. Other categories include false-positive triage ($12.9\%$), retrieval miss ($7.9\%$), and a residual long tail.

\paragraphhead{Reformulation (RQ-A3, not measurable in v1).} The reformulator's gate opens on $179$ of $1{,}008$ windows ($17.8\%$), and the primitive emits a reformulated query under its bounded action vocabulary. The primitive's Hit@K contribution cannot be measured on this paper's predictions-backed retrievers, which key into a cached prediction by window identifier rather than by query text and therefore do not respond to text-level reformulation. The framework is in place; a live-retrieval re-run is required to measure the lift and is left to future work.
